% Library experiments. Compiled by library_notebook.tex, not by main.tex.
%
% A separate record from experiments.tex because this is a different arc. That
% file is the AME benchmark programme -- E1-E22, selection rules, oracles,
% run-to-run variance. This one is about whether designs can be cut into
% fragments and reassembled, which is the question a combinatorial library
% actually turns on, and which none of E1-E22 was built to answer.
%
% Same discipline: pre-registrations stay next to the results that confirmed or
% killed them, and retracted claims stay visible rather than being quietly
% rewritten. Three claims from the AME arc were overturned in the week before
% this file was started, and they are recorded here because the overturning is
% what justifies the change of thesis.

\section{Library experiments: fragments, recombination, and modularity}

\subsection*{Why a separate record}

The AME programme asked whether reinforcement learning produces \emph{better}
sequences. This arc asks whether it produces \emph{reusable} ones. A
combinatorial library is assembled from parts: $P$ passing parents cut into $k$
slots yield $P^k$ constructs from $k \cdot P$ synthesised fragments, and that
product-versus-sum asymmetry is the only reason to build libraries this way. But
$P^k$ counts constructs, not working constructs. If the recombinants do not
fold, the number is fiction.

Nothing in the AME slate measured that, and nothing in the published literature
we are aware of does either.

\subsection*{Three claims this arc retired}

\plan{Recorded because the change of thesis rests on them, and because each was
asserted here before the data that killed it.

\textbf{The diversity/fidelity frontier.} On one backbone at a 96-fold budget
the policy dominated the entire temperature sweep. Measured across 13 backbones
at a \emph{matched} budget, it wins both axes on 6/13 and loses on the easy end
where sampling already works. The original comparison excluded the policy's
2{,}400-fold training debt.

\textbf{Rescue.} ``Sampling yields nothing on hard backbones; the policy rescues
them'' rested on $0/10{,}000$ at $T=1.0$ alone. Given the same budget spread
across the temperature sweep, the baseline finds 116 and 1 passing designs on
those backbones. Clean rescue is 1/13, with a counterexample beside it (M0078:
policy 0, baseline 11).

\textbf{Difficulty axes are uncorrelated.} Claimed from three backbones, where
the hard one outperformed the medium one. Across 13, policy pass rate correlates
with the published dead-backbone fraction at $r = -0.776$. Both methods find the
same backbones hard.}

\subsection*{L1 --- Does recombination work at all?}

\plan{Pre-registered. Split passing designs at $k=4$ equal boundaries,
recombine, exclude exact parents, fold best-of-5 under AlphaFold3. Retention is
the chimera pass rate; parents are passing by construction, so no ratio is
involved. 256 chimeras per cell, 21 cells across 13 backbones.}

Retention spans \SIrange{0.4}{95}{\percent}. On the easier half it is high
enough that the combinatorial multiplier is real: a few hundred parents give
$\sim$\num{e9} constructs of which most fold.

\subsection*{L2 result --- RL-designed fragments recombine better}

\begin{center}
\begin{tabular}{rlrr}
\toprule
dead\% & backbone & sampled & RL-trained \\
\midrule
16 & M0664\_2dhn & 90.6 & \textbf{94.9} \\
55 & M0097\_1ctt & 59.8 & \textbf{72.7} \\
77 & M0255\_1mg5 & 25.8 & \textbf{79.3} \\
85 & M0315\_1ey3 & 49.6 & \textbf{78.1} \\
86 & M0500\_1e3i & 10.9 & \textbf{53.5} \\
95 & M0058\_1cju & 12.9 & \textbf{44.9} \\
97 & M0050\_1dbt & \textbf{11.7} & 5.5 \\
97 & M0904\_1qgx & 12.9 & \textbf{47.3} \\
99 & M0365\_1pfk & 0.4 & \textbf{5.5} \\
\bottomrule
\end{tabular}
\end{center}

\textbf{8 of 9, mean paired advantage \SI{+23.0}{points}}; mean retention
\SI{48.1}{\percent} against \SI{28.5}{\percent}. Nothing in the objective asked
for this --- the policy was trained on motif RMSD alone.

Two things make it load-bearing rather than incidental. It holds \emph{where the
policy loses on parent count}: on M0664 the baseline has \num{2273} passing
parents to the policy's 243, nearly \num{10}$\times$, and its chimeras still
pass less often. And fragment-space \emph{size} at matched parent count is
exactly $1.00\times$ between arms --- at 45 residues no two passing designs
share a fragment, so library size is determined entirely by parent count.
Retention is the only axis on which the arms differ.

\subsection*{L3 result --- the surviving parts pool is also more diverse}

The diversity that matters for a library is over fragments that still work after
shuffling, not over all designs.

\begin{center}
\begin{tabular}{lrrrr}
\toprule
 & \multicolumn{2}{c}{survival fraction} & \multicolumn{2}{c}{$U$ over survivors (rarefied)} \\
 & sampled & RL & sampled & RL \\
\midrule
mean & 0.52 & \textbf{0.84} & --- & --- \\
backbones won & \multicolumn{2}{c}{9/9 to RL} & \multicolumn{2}{c}{8/8 to RL} \\
\bottomrule
\end{tabular}
\end{center}

\plan{\textbf{This reverses an earlier finding and the reversal matters.} Over
all passing \emph{designs}, $T=1.0$ carries more substitutions than the policy
(\num{1988} against \num{1045}). Over \emph{usable fragments}, rarefied to a
matched survivor count so pool size cannot explain it, the policy carries more
on every backbone. Both are true. The first was measuring the wrong object for a
library built from parts, and the diversity claim should be stated over
survivors or not at all.}

\subsection*{L4 --- Training for recombinability}

\plan{Pre-registered, running. Fragment-level GRPO: sample 8 parents and never
fold them; build 32 chimeras by rotation so every fragment appears in exactly 4;
fold those; give each fragment an advantage equal to the mean motif RMSD over
its chimeras, standardised within slot. The $[B, L]$ advantage is consumed by
the existing per-token surrogate without changing \texttt{grpo.py}.

Design choices were simulated against known fragment effects rather than
guessed: rotation beats random permutation at this size; simple averaging beats
least-squares regression at 32 chimeras (32 observations, 32 parameters); and
scoring on continuous RMSD avoids the tied advantages a four-chimera binary
would produce. Recovery of the true pairwise fragment ordering is
$\approx$\SI{80}{\percent}, falling to $\approx$\SI{67}{\percent} once fragment
interactions are included --- diluted, but enough to shift a distribution.

\textbf{Prediction.} Retention and survival fraction rise above the policy arm
on all three backbones (M0097, M0255, M0315). A small effect is consistent with
the signal quality above and is not by itself evidence against the idea.

\textbf{Kill criterion.} Fragment collapse. The objective can be satisfied by
emitting near-identical sequences: every chimera then reproduces a parent and
passes. That shows as survival fraction climbing toward 1.0 while $U$ over
survivors falls. Retention alone cannot distinguish it; the pair can. If it
triggers, the term is gated on chimeras differing from parents, \emph{not}
reweighted --- the weight pilot measured that adding a diversity reward at any
weight from 1 to 4 drove pass rate to 0/256 with median motif RMSD
\SIrange{5.9}{11.8}{\angstrom}.}

\subsection*{L5 result --- the usable granularity, and where it ends}

\plan{Pre-registered: retention falls monotonically with $k$; the RL arm's
advantage persists across $k$; if it vanishes by $k=8$ the claim is ``designs
tolerate quarter-swaps'', which is much narrower than ``designs are modular''.}

42 cells: three backbones $\times$ two arms $\times$
$k \in \{2,4,6,8,12,16,24,36\}$, 256 recombinants each,
\numproduct{10752} designs folded best-of-5.

\begin{center}
\begin{tabular}{rrrrcrrr}
\toprule
& \multicolumn{3}{c}{sampled} & & \multicolumn{3}{c}{RL-trained} \\
\cmidrule(r){2-4}\cmidrule(l){6-8}
$k$ & M0097 & M0255 & M0315 & & M0097 & M0255 & M0315 \\
\midrule
2  & 64.1 & 35.9 & 56.6 & & 81.2 & 80.5 & 86.3 \\
4  & 59.8 & 25.8 & 49.6 & & 72.7 & 79.3 & 78.1 \\
6  & 50.4 & 27.0 & 46.9 & & 77.3 & 82.4 & 80.1 \\
8  & 51.6 & 25.8 & 43.4 & & 73.4 & 77.0 & 77.3 \\
12 & 44.9 & 21.9 & 42.6 & & 71.1 & 68.0 & 73.8 \\
16 & 39.8 & 19.5 & 37.9 & & 57.4 & 61.3 & 74.2 \\
24 & 28.1 & 12.5 & 22.3 & & 48.4 & 36.3 & 64.8 \\
36 & 12.1 &  9.8 &  7.8 & & 25.8 & 15.2 & 42.2 \\
\bottomrule
\end{tabular}
\end{center}

\textbf{The prediction was half right.} Retention does fall, but not
monotonically and not early: it is flat through $k \approx 8$ --- 22-residue
fragments still recombine --- and only then declines. The advantage does
persist, and widens.

Taking \SI{50}{\percent} as the point where a library stops being worth
assembling (two constructs folded per usable one), the largest workable $k$ is

\begin{center}
\begin{tabular}{lccc}
\toprule
 & M0097 & M0255 & M0315 \\
\midrule
sampled     & 8  & $<2$ & 2  \\
RL-trained  & \textbf{16} & \textbf{16} & \textbf{24} \\
\bottomrule
\end{tabular}
\end{center}

\textbf{Training extends the usable granularity by roughly $2$--$8\times$ in
$k$.} Library size goes as $\prod_i |F_i|$, so that is the largest practical
effect measured in this project --- far larger than the difference in parent
counts, which runs the other way on several backbones.

\plan{\textbf{Fragment collision is a second degradation signature, and it needs
no folding.} Through $k=6$ every parent contributes a unique fragment to every
slot. Past that, designs start sharing fragments: M0097's sampled arm falls to
\SI{64}{\percent} of parents distinct per slot at $k=12$ and \SI{16}{\percent}
at $k=36$, while the RL arms stay above \SI{80}{\percent} to $k=16$. Those
\num{1042} sampled parents are far more redundant than the count suggests, which
partly explains why they recombine worse despite outnumbering the policy's 178
by six to one.}

\plan{\textbf{Library size is deliberately not plotted against $k$.} At a fixed
synthesis budget both arms build the same-size library --- each slot takes
budget$/k$ fragments regardless of which arm produced them --- so those curves
lie on top of one another. That is itself the finding: the arms do not differ in
how large a library they can build, only in what fraction of it folds. An
earlier version of the figure used $P^k$, which overstates the library by
$10^{15}$ at $k=36$ once fragments collide; it now uses the product of measured
distinct fragments.}

\subsection*{L6 result --- boundary placement does not matter}

\plan{Pre-registered in \texttt{results/fragments.tex} before any of this ran,
and quoted unchanged: \emph{``If they do not, the structural rationale is
decorative and should be dropped.''}}

Uniform boundaries against boundaries moved to the nearest coil residue (within
12 positions, secondary structure from \texttt{biotite.annotate\_sse}), $k=4$,
seven backbones $\times$ two arms, 256 recombinants per cell.

\begin{center}
\begin{tabular}{lrr}
\toprule
 & mean $\Delta$ & \\
\midrule
all 14 cells   & \SI{+0.4}{points} & 95\,\% CI $[-1.3, +2.0]$, $t=0.42$, $p=0.68$ \\
sampled arm    & \SI{-0.7}{points} & $n=7$ \\
RL arm         & \SI{+1.5}{points} & $n=7$ \\
\bottomrule
\end{tabular}
\end{center}

Every individual cell sits within $1.6$ standard errors of zero. By its own
pre-registered criterion the structural rationale is dropped: \textbf{equal
slots throughout}.

\plan{\textbf{A sub-claim retracted.} On the first six cells the RL arm appeared
to benefit more ($+5.9$, $+0.8$, $+1.6$), and that was written up here as
``directionally consistent with RL designs having more coherent secondary
structure''. With seven cells per arm the RL spread is $-5.1$ to $+5.9$ and the
means are $+1.5$ against $-0.7$. It was pattern-reading across six numbers all
within one or two standard errors of zero, and it is withdrawn.}

\plan{\textbf{Equal slots are also easier to assemble, but the margin is
smaller than first claimed.} Structural boundaries give fragment lengths
spanning \SIrange{108}{162}{bp}, a $1.28$--$1.50\times$ spread; since moles
scale as $1/\text{length}$ at fixed mass, equimolar Golden Gate pooling then
needs per-part normalisation. That spread is a tunable consequence of the
$\pm 12$ residue shift window, not a property of structural boundaries, and
$1.3\times$ is workable in practice. So the honest statement is that structural
boundaries buy no measurable retention and cost a little length uniformity ---
not that they are unusable.}

\subsection*{What the two measured results mean together}

\plan{The L2/L3 result and the L5 result are different \emph{kinds} of evidence
and should not be presented as equally strong.

The panel policies were trained on parent motif RMSD alone --- chimeras appear
nowhere in their objective. Their retention advantage is therefore a property
measured on a quantity nothing in training optimised, which is the strongest
form this claim can take.

The L4 arm trains \emph{on} chimera quality, so its retention is the same kind
of quantity it optimised, under a different predictor. The cross-oracle swap is
the only thing between the training signal and the reported number. That is
unavoidable when training for the thing you want to measure, but it means the L4
result is a comparison against the panel arm, not a standalone number.}
